% Nouveau EGA V (cote 25): archive cote 25, site page 49 (PDF page 50).
% typescript. Draft compared: none.
% AI-transcribed by Claude Opus 5.5 (claude-opus-5-5, adaptive thinking, effort high), Message Batches API, batch ega25b (msgbatch_01A7n5YuXubUPNK1zsepi8RG), 2026-11-13; original AI transcription, subsequently checked against both supplied scan parts; see PAGES_CHECKED.tsv and CHECK_DECISIONS.tsv.
\scanpage{25/49}
% NOTE: The small mark after dimension at most2 has unsettled character and layer. The first cancelled component index retains i/a uncertainty. The final des conditions has a faint displaced type shadow.

\mspage{\hand{45}}

toute composante irréductible $X_i$ de dimension 1  de X soit 
formée de points lisses de X et soit  réglée relativement à f, et que
pour tout point fermé  x de X tel que $\dim_x X \hand{\geqslant} 2$, on ait $\mathrm{prof}_x X \hand{\geqslant} 2$,
(conditions automatiquement vérifiées si, X est géom normale et toutes ses
composantes irréductibles sont de dim$\hand{\geqslant}$ 2).

b) Supposons X géom. normale ,  que l'ensemble  des $\hand{\xi \in}\, \overset{\hand{\vee}}{P}$
tels que $\phi_{\hand{\xi}}$ soit $X_{k(\hand{\xi})}$\hand{-}régulier et $Y_{\hand{\xi}}$  géom normale ait un complé-%
mentaire de codimension $\hand{\geqslant}$2, il f et s que  toute composante
irréductible $X_i$ de X de dimension$\hand{\leqslant}$2\unsure{,|x} soit formée de point lisses de
X et soit réglée rel à f, et que de plus pour tout point fermé x de X
tel que $\dim_x X \hand{\geqslant} 3$, on ait $\mathrm{prof}_x X \hand{\geqslant}$ 3.





\uline{Remarque} 8.15. Dans 8.6. , 8.13. et 8.14. nous avons fait sur X
l'hypothèse $(S_k)$ (resp. $(R_k)$, resp: séparable, resp. géom normal ) que
nous désirions receuillir comme conclusion pour les sections hyperplanes
sauf pour $\xi$ dans un ensemble exceptionnel de codim\unsure{e|m}nsion $\hand{\geqslant}$ 2. Cela ne restreint
pas la généralité (et il n'aurait pas coûté plus cher à vrai dire
de se dispenser de cette hypothèse préalable), car on voit aussitôt
à l'aide des résultats de par. 3.4. et 5.12. que si X ne satisfait pas à
l'hypothèse en question, donc (en vertu de par.5) s'il existe un point
\uline{fermé} x en lequel cette hypothèse est en défaut, alors pour tout $\hand{\xi}$ tel
que $\phi_{\hand{\xi}}$ soit $X_{k(\hand{\xi})}$-régulier (condition n'éliminant qu'un ens de codim$\hand{\geqslant}$2)
et tel que $x \hand{\in} Y_{\hand{\xi}}$ (condition qui décrit un ensemble de codimension 1
exactement), $Y_{\hand{\xi}}$ ne satisfait pas non plus à ladite hypothèse en x ,
donc l'ensemble exceptionnel $Z \hand{\subset} \overset{\hand{\vee}}{P}$ est de codimension 1 et non 2.
(J'ai peut-être charrié un peu dans le cas $(R_k)$, où il faut encore
des conditions $(S_1)$ et d'équidimensionnalité peut-être.....).

\placedmargin{\hand{49}}
